\section{Selección de BindsNET como framework de desarrollo}

La selección de BindsNET como framework de desarrollo se fundamentó en una evaluación comparativa de las opciones disponibles para la simulación de SNNs orientadas al aprendizaje automático \cite{Hazan2018_BindsNET}. BindsNET proporciona objetos y métodos de software que permiten la simulación de grupos de diferentes tipos de neuronas (bindsnet.network.nodes), así como de diferentes tipos de conexiones entre ellas (bindsnet.network.topology). Estos pueden combinarse en un único objeto bindsnet.network.Network, responsable de coordinar la lógica de simulación de todos los componentes subyacentes \cite{bindsnet_docs}.

Las características técnicas que determinaron esta selección incluyen su integración nativa con PyTorch (herramienta con la que se cuenta experiencia previa), lo que facilita la implementación en plataformas computacionales de alto rendimiento tanto en CPU como GPU. Adicionalmente, BindsNET ofrece una sintaxis concisa y orientada al usuario que resulta particularmente adecuada para prototipado rápido, aspecto crucial durante las fases iniciales de desarrollo algorítmico \cite{bindsnet_docs}.

Los detalles sobre la instalación y configuración del entorno de desarrollo, incluyendo la resolución de conflictos de dependencias y problemas de compatibilidad, se encuentran documentados en el Anexo~\ref{anexo:manual-instalacion}.